\section{Outlook}
\label{sec:outlook}

\noindent\emph{Status: draft -- the milestone list and the acceptance criteria are transcribed from the roadmap
that governs the work; the open problems are stated as problems, and none of the milestones after M0 is claimed
to be working.}

What remains is not a refinement of what is finished. The cavity solver of Sections~\ref{sec:eigen}
and~\ref{sec:results} is accurate and its energy is exact, but the package cannot yet radiate: there is no
absorbing boundary, no near-to-far-field transform and no port, so every problem it can address is enclosed by
walls. The roadmap is therefore a list of capabilities rather than a list of improvements, and each entry carries
an acceptance test expressed in terms of a quantity the code cannot currently compute. That is deliberate: a
milestone is done when its number is reproduced by the test suite, not when the module exists.

\begin{table}[H]
  \centering
  \caption{The roadmap. Status, deliverable and acceptance criterion for each milestone, in dependency order
  rather than order of value. Only M0 is complete; M1 and M5 are in progress.}
  \label{tab:milestones}
  \begin{tabular}{@{}lp{1.8cm}p{5.6cm}p{5.6cm}@{}}
    \toprule
    & Status & Deliverable & Acceptance criterion \\
    \midrule
    M0 & done & time-domain core with the exact invariant; dense cavity eigensolver & thirty tests pass; box
    modes within $1.3$\,\% ($1.213$\,\% coarse, $0.615$\,\% fine) and second-order ratio $0.507$ \\
    M1 & in progress & first-order Mur absorbing condition on the six outer walls, then a stretched-coordinate
    PML layer & PEC sphere of radius $0.1$ m: bistatic RCS within $0.5$ dB of the Mie series at $ka=1$
    (\SI{477}{\mega\hertz}) and $ka=3$ \\
    M2 & planned & near-to-far-field transformation over a Huygens surface, one transient run transformed &
    half-wave dipole at \SI{300}{\mega\hertz}: $\operatorname{Re}Z_{\mathrm{in}}=73\pm5$\,\si{\ohm},
    $|\operatorname{Im}Z_{\mathrm{in}}|<20$\,\si{\ohm}, $2.15\pm0.15$\,dBi \\
    M3 & planned & two-dimensional modal port eigenproblem, mode overlaps, de-embedding & WR-90 waveguide,
    TE$_{10}$ at \SI{6.557}{\giga\hertz}: $|S_{11}|<-30$\,dB, $|S_{21}|=1.000\pm0.001$, phase within
    \SI{1}{\degree} \\
    M4 & planned & KiCad s-expression reader, staircase voxelisation, stack-up to material parameters & microstrip
    $Z_0$ within $5$\,\% of Hammerstad--Jensen; patch resonance within $3$\,\% of $c/(2L\sqrt{\eps_{\mathrm{eff}}})$ \\
    M5 & in progress & gradient null-space deflation followed by a matrix-free eigensolver beyond the dense
    threshold & $10^5$ free edges, six modes in under five minutes, within $0.1$\,\% of a dense solve \\
    M6 & planned & $\mathbf{A}$--$V$ magnetoquasistatics with eddy currents; parareal time integration & aluminium
    slab, skin depth $2.7$ mm at \SI{1}{\kilo\hertz} and $11.0$ mm at \SI{60}{\hertz}, within $2$\,\% \\
    \bottomrule
  \end{tabular}
\end{table}

Two of these entries are gates rather than features. M1 is the gate to every open-region question: without an
absorbing boundary there is nothing to compare with a Mie series, and without a far-field transform there is no
radiation pattern to draw, so M2 has no meaning until M1 has produced a field that leaves the domain. M5 is the
gate to every engineering-sized geometry: the dense solver of Section~\ref{sec:eigen} refuses more than $3000$
free edges, which is a cavity of a few wavelengths and no more, and nothing about a printed antenna fits inside
that budget. The other four are independent of each other, which is why they can be reordered as time and
interest allow.

\subsection{Open problems}

The following are difficulties we do not yet know how to price, not features we have promised. Each is stated as
the question we would have to answer before the corresponding milestone could be called finished.

\begin{enumerate}[leftmargin=1.6em,itemsep=3pt]
  \item \textbf{Whether the dispersion is worth repairing.} At about four cells per wavelength the tracked front
        travels at $0.904\,c_0$, and \eqref{eq:phase} says why. The open question is what to do about it: raise
        the order of the spatial operator, refine until the bias is negligible, or accept it and quote
        frequencies with a known sign of error. Higher order means giving up the two-point stencil on which the
        exact identity \eqref{eq:invariant} rests, so the cheap repair and the strong conservation result are in
        direct competition, and we do not know which we would choose.
  \item \textbf{The staircase error on curved geometry.} Voxelisation of a curved surface produces a boundary
        off by up to half a cell, and the resulting frequency error is not second order. A patch antenna imported
        from a layout file will be a staircase, and by how much that degrades the resonance is unknown to us; it
        may be that a fine local grid is enough, or it may be that a conformal treatment of the boundary is
        unavoidable and the structured-grid advantage is lost exactly where it is most wanted.
  \item \textbf{Deflating the null space at scale.} The gradient modes that make \eqref{eq:geneig} singular
        number in the thousands on the grids that matter, and the projector that removes them is a dense-ish
        operator applied inside an iterative solve. Whether the deflated iteration converges in a time an engineer
        will wait for, and whether it holds the $0.1$\,\% accuracy the dense path gives, is the substance of M5;
        the shift-invert route was tried and rejected, which tells us the obvious approach is wrong but not that
        the replacement is right.
  \item \textbf{The global time step.} The bound \eqref{eq:cfl} is set by the smallest cell anywhere, so one
        narrow gap or one thin conductor dictates the step count for the whole domain, and a graded grid can
        multiply the number of steps by four orders of magnitude. Local time stepping is the textbook answer and
        it interacts badly with an exact global invariant, since a scheme with different steps in different
        regions does not have a single conserved quadratic form of the simple sort.
  \item \textbf{Where the boundary goes.} Mur's condition and a Berenger layer both assume an outgoing wave and
        both are approximate on a staggered grid, and the reflection they leave behind sets the floor on any
        scattering result. Absorbing more than $20$ dB, and then more than $40$ dB, is a numerical question we
        have not answered; and a PML needs a diagonal tensor material in the layer, which the present scalar
        material layer cannot express, so M1 begins with an extension of the data structure rather than with an
        algorithm.
\end{enumerate}

\subsection{Scope}

We want to be plain about what the report as a whole claims. \texttt{emfit} is, today, a verified solver for
closed, isotropic, lossless problems on structured grids: it conserves the discrete energy to rounding, it
reproduces the analytic spectrum of a conducting box to about one per cent on coarse grids with second-order
convergence, and it propagates a wavefront at the speed its own dispersion relation predicts. It is not yet a tool
for the questions posed in Section~\ref{sec:theory} about antennas, laminate or scatterers, because those need a
field that leaves the domain and a port through which power is measured, and neither exists. The milestones above
are acceptance tests for that transition; the open problems are the parts of it we cannot yet cost. Progress will
be recorded here as it happens, including the parts that go badly.
